\chapter*{Notation and Conventions}
\addcontentsline{toc}{chapter}{Notation and Conventions}
\markboth{Notation and Conventions}{Notation and Conventions}

This monograph presents a completely self-contained, exact, and constructive formulation of $1+1$-dimensional lattice bosonization within algebraic quantum field theory (AQFT). Unlike traditional continuum heuristic treatments that rely on divergent operator products, formal distribution limits, or unbounded operators, all mathematical statements established here are exact theorems on strictly finite-dimensional carrier spaces, finite lattices, or explicitly graded polynomial algebras.

To maintain clarity and consistency across all chapters, this section establishes the unified mathematical notation used throughout the text.

\section*{Lattice Geometry and Discrete Kinematics}

\begin{itemize}
    \item \textbf{Lattice size:} $L \in \mathbb{N}$ denotes the number of discrete spatial lattice sites, strictly positive ($L \ge 1$). Throughout the fermionic and bosonization development, $L$ is assumed to be an even integer, parametrized by a half-length $h \ge 1$ such that $L = 2h$.
    \item \textbf{Spatial lattice:} $\Lambda = \mathbb{Z}/L\mathbb{Z}$ is the discrete 1D periodic chain of sites $x, y \in \Lambda$, endowed with the standard quotient group structure modeling periodic boundary conditions.
    \item \textbf{Momentum band:} $\Lambda^*$ denotes the centered discrete Brillouin zone of $L$ integers:
    \[
    \Lambda^* = \left\{ k \in \mathbb{Z} \;\middle|\; -L < 2k \le L \right\}.
    \]
    When $L = 2h$, this interval is explicitly $\Lambda^* = \{-h+1, -h+2, \dots, h-1, h\}$. This choice retains the positive Nyquist mode $k = +h = L/2$ and excludes $-h$, ensuring an exact bijection with $\Lambda$ while deliberately breaking reflection symmetry.
    \item \textbf{Projection and Section:}
    \begin{itemize}
        \item $\pi : \mathbb{Z} \to \Lambda$ denotes the canonical quotient projection $\pi(k) = k \bmod L$.
        \item $r : \Lambda \to \Lambda^*$ denotes the unique centered section of $\pi$ landing in $\Lambda^*$.
    \end{itemize}
    \item \textbf{Transported band arithmetic:}
    \[
    k \oplus p = r(\pi(k) + \pi(p)), \qquad \ominus k = r(-\pi(k)).
    \]
    The discrepancy between integer addition $k+p$ and band addition $k \oplus p$ is governed by the Umklapp integer $w \in \{-1, 0, 1\}$ such that $k+p = (k \oplus p) + wL$.
\end{itemize}

\section*{Difference Calculus and Fourier Analysis}

\begin{itemize}
    \item \textbf{Shift and differences:} For functions $f : \Lambda \to \mathbb{C}$, the shift operator is $(E f)(x) = f(x+1)$, the forward difference is $\Delta = E - I$, and the backward difference is $\nabla = I - E^{-1}$.
    \item \textbf{Fourier transform:} For $k \in \Lambda^*$ and $x \in \Lambda$, the character is $\chi(k,x) = \exp(2\pi i k x / L)$. The discrete Fourier transform $\mathcal{F}$ maps spatial functions to momentum space unitarily, satisfying exact Parseval and Plancherel identities over the finite cyclic group.
\end{itemize}

\section*{Fermionic Fock Spaces and Observable Nets}

\begin{itemize}
    \item \textbf{CAR Algebra:} The canonical anticommutation relations on the momentum band $\Lambda^*$ are generated by operators $c_k, c_k^\dagger$ satisfying:
    \[
    \{c_k, c_p^\dagger\} = \delta_{k,p} I, \qquad \{c_k, c_p\} = 0, \qquad \{c_k^\dagger, c_p^\dagger\} = 0, \quad \forall k, p \in \Lambda^*.
    \]
    \item \textbf{Spatial fermions:} Constructed via the finite discrete Fourier transform:
    \[
    \psi_x = \frac{1}{\sqrt{L}} \sum_{k \in \Lambda^*} \chi(-k, x)\, c_k, \qquad \psi_x^\dagger = \frac{1}{\sqrt{L}} \sum_{k \in \Lambda^*} \chi(k, x)\, c_k^\dagger.
    \]
    \item \textbf{Local nets:} For an interval $I \subset \Lambda$, $\mathfrak{A}(I)$ is the von Neumann / $C^*$-subalgebra generated by $\{\psi_x, \psi_x^\dagger \mid x \in I\}$.
\end{itemize}

\section*{Vacuum Structure, Charge Sectors, and Energy Budgets}

\begin{itemize}
    \item \textbf{Dirac sea vacuum:} $|\Omega\rangle$ denotes the half-filled reference sea configuration $S_0 = \{-h+1, \dots, 0\}$.
    \item \textbf{Admissible charge sectors:} For $N \in \mathbb{Z}$ with $-h \le N \le h$, the sector ground configuration is $S_N = \{-h+1, \dots, N\}$, with ground state $|N\rangle_0$ of total particle number $h + N$.
    \item \textbf{Excitation energy:} For an occupation configuration $\sigma \subset \Lambda^*$ of cardinality $h+N$, the excitation energy above the sector ground is:
    \[
    E(\sigma) = \sum_{k \in \sigma} k - \sum_{k \in S_N} k \ge 0.
    \]
    \item \textbf{Energy budget subspace:} For sector $N$ and energy cutoff $K \ge 0$:
    \[
    B(N, K) = \bigoplus_{E=0}^K \mathcal{H}(N, E),
    \]
    where $\mathcal{H}(N,E)$ is the subspace spanned by configurations with charge $N$ and excitation energy $E$.
\end{itemize}

\section*{Density Modes and Bosonic Representation}

\begin{itemize}
    \item \textbf{Normal-ordered density modes:} For mode integer $m \neq 0$:
    \[
    :\!\rho_m\!: = \sum_{k \in \Lambda^*} : \! c_{k \ominus m}^\dagger c_k \!:
    \]
    \item \textbf{Heisenberg / Kac--Moody algebra:} Under explicit margin conditions ($2M + K + |N| \le h$), the density modes satisfy the exact commutation relations on $B(N,K)$:
    \[
    \left[ :\!\rho_m\!:, :\!\rho_{-n}\!: \right] = m\, \delta_{m,n}\, I.
    \]
    \item \textbf{Bosonic partition states:} For an integer partition $\lambda = (\lambda_1 \ge \lambda_2 \ge \dots \ge \lambda_\ell > 0)$ of size $|\lambda| = \sum \lambda_i = E$:
    \[
    |\lambda; N\rangle = :\! \rho_{\lambda_1} \!: :\! \rho_{\lambda_2} \!: \dots :\! \rho_{\lambda_\ell} \!: |N\rangle_0.
    \]
    \item \textbf{Haldane completeness:} The states $\{|\lambda; N\rangle \mid |\lambda| = E\}$ form an exact basis of $\mathcal{H}(N, E)$, and $\bigcup_{E=0}^K \{|\lambda; N\rangle \mid |\lambda| = E\}$ forms a basis of $B(N, K)$.
\end{itemize}

\section*{Conventions on Constructive Realization}

Every mathematical theorem in this monograph corresponds to a fully verified, kernel-checked theorem in the Lean 4 formalization. Rather than polluting the mathematical discourse with code identifiers, tactic transcripts, or backticks, the text presents rigorous, complete human-readable proofs. 

Where constructive methods deviate from traditional non-constructive continuum habits (such as the choice of asymmetric Nyquist modes, finite decidable filters, quotient injectivity, or polynomial state carriers), the text highlights the underlying mathematical rationale in dedicated \emph{Formalization Note} remarks.
